\chapter{面向链式推理的数据流与调度优化}
\label{chap:dataflow}

LASA 提供了层级可配置的硬件执行基础，但真实 YOLOv3-tiny 链式推理的端到端延迟仍不能只由阵列峰值吞吐决定。权重加载、IFM 打包、DMA 启动、tile 边界同步、partial-PSUM 反馈、pass 准备和 OFM 输出都会形成非计算开销。因此，本文将优化目标写为
\begin{equation}
T_{noncomp}=T_{layer}-T_{compute\_fire}.
\end{equation}
其中 $T_{compute\_fire}$ 表示阵列真正接受输出 pixel 计算的周期数，其余部分均属于等待、搬运、反馈或边界开销。

本章首先提出 KCS 三维分块执行模型，将不同层统一映射为 K-pass、Cout block 和 spatial tile。随后围绕三类非计算瓶颈提出三类优化策略：面向 PS/PL 交互的批量流式协同策略 BSD（Batched Streaming Dataflow）、面向 IFM 布局转换的片上重排与重放策略 OCRR（On-Chip Reorder and Replay），以及面向 K-pass 边界气泡的反馈重叠与预取策略 OPF-P（Overlapped PSUM Feedback and Prefetch）。

\begin{figure}[htb]
\centering
\includegraphics[width=\textwidth]{\detokenize{img/picture/时序优化策略.png}}
\caption{BSD、OCRR 与 OPF-P 三类优化策略总览}
\label{fig:strategy_overview}
\end{figure}

\section{KCS 三维分块执行模型}

KCS 执行模型由 K 维、输出通道和空间三个调度维度组成。K 维分块由 ROWS 决定，最终配置 ROWS=18，因此 $K_{tile}=18$。输出通道分块由 COLS 和 dual-lane INT8 column 决定，每列产生两个输出通道 lane，因此
\begin{equation}
Cout_{tile}=2\times COLS.
\end{equation}
最终配置 COLS=8，对应 $Cout_{tile}=16$。空间分块由 tile 的输出行范围决定，用于限制片上 IFM/OFM buffer 和单次 AXI DMA 长度。

图~\ref{fig:tiling_schedule} 给出三维调度关系。每个 spatial tile 内，调度器先遍历输出通道块，再对当前输出通道块执行所有 K pass。首个 K pass 从 bias 初始化 partial sum；非首 pass 读取上一轮 partial-PSUM；final pass 将结果送入 requant、activation、pooling 和 OFM path。

\begin{figure}[htb]
\centering
\includegraphics[width=\textwidth]{\detokenize{img/picture/KCS分块策略.png}}
\caption{KCS 三维分块执行模型}
\label{fig:tiling_schedule}
\end{figure}

KCS 不是单独的优化策略，而是 LASA 上的统一映射抽象。它回答不同 CNN 层如何映射到固定阵列规模，而 BSD、OCRR 和 OPF-P 则分别针对映射后暴露出的 PS/PL 交互、数据排布和 pass 边界开销进行优化。

\section{BSD：面向 PS/PL 交互的批量流式协同}

BSD 的目标是降低 PS 逐请求服务和 DMA 频繁启动造成的系统协同开销。早期 per-request 模式中，硬件每需要一个 bias block、weight tile 或 IFM line 时向软件发出请求，软件随后打包数据、启动 DMA 并等待完成。这种方式适合小型 smoke test，但在 Conv6 等真实层中会产生大量服务事件，A53 侧开销和 DMA 启动开销会淹没阵列计算。

BSD 将每个 spatial tile 所需的 bias、weight 和 IFM 数据预先打包成连续 AXI stream。软件在 tile 开始时分别启动三路输入 DMA，硬件通过 packet count 恢复内部边界，只在整个 stream 末尾依赖 TLAST。这样，DMA 启动次数从每个请求一次降低为每个 tile 每类输入一次，AXI backpressure 也由 PL 自然吸收。

权重数据在模型中通常按输出通道、输入通道和 kernel 维度组织，而硬件消费顺序是 $Cout block \rightarrow K pass \rightarrow lane \rightarrow channel$。本文在生成 Vitis header 时直接产生硬件消费顺序的权重数组，使 weight DMA 可以读取 ELF 只读段中的连续数组，不再需要运行时 scratch copy 或 repack。PL 端 weight loader 使用 64-bit AXI beat 写入多个 byte bank，进一步降低 weight load 等待。预打包权重虽然作用于数据布局，但其主要收益是减少运行时 PS 服务和 DMA 准备成本，因此归入 BSD。

\section{OCRR：面向 IFM 布局转换的片上重排与重放}

OCRR 的目标是缓解软件自然 HWC 布局与 systolic array K-lane 消费顺序之间的不匹配。IFM 输入存在两种路径。第一种是 prepacked IFM stream：软件提前把每个 K pass 需要的 IFM lane 按硬件顺序打包。该方式 PL 简单，但 A53 侧打包成本较高。第二种是 raw-HWC stream：软件发送自然 HWC tile，PL 在片上缓存后按 pass 和窗口顺序 replay。

图~\ref{fig:strategy_overview} 中间部分给出 OCRR 的数据排布思路。prepacked 路径把重排压力放在软件；raw-HWC 路径把重排压力转移到 PL cache/replay，但可减少 A53 逐 pass 复制和索引开销。对 3x3 层，cache materializes 每个输出 pixel 的 3x3 窗口，并按 18 lane 输出给阵列；对 1x1 层，cache 按连续通道向量 replay。

OCRR 还包括后端 full-tile 调度。Conv5、Conv6 和 Conv8 的空间尺寸均为 $13\times13$。早期调度将这些层拆成 $4+4+4+1$ 行 spatial tile，可以降低片上缓存容量需求，但每个 tile 都需要重复启动 bias、weight、IFM 和 OFM 流程。随着 raw-HWC cache 使用 URAM 扩容，后端层可以使用完整 $13\times13$ tile，从而减少 tile 边界固定成本。容量条件可近似写为
\begin{equation}
tile\_pixels \times \left\lceil\frac{C_{in}}{2}\right\rceil \le HWC\_CACHE\_DEPTH.
\end{equation}
在最终 full-tile build 中，Conv5/8 需要 $13\times13\times\lceil256/2\rceil=21632$ 个 materialized word，Conv6 需要 $13\times13\times\lceil512/2\rceil=43264$ 个 materialized word。raw-HWC 并非对所有层都收益为正；Conv3 raw-HWC 虽然可功能通过，但由于 tile/FIFO 容量和 replay 开销，板级结果不优于推荐配置，因此作为边界分析而非默认路径。

\section{OPF-P：面向 K-pass 边界的反馈重叠与预取}

OPF-P 的目标是隐藏多 K pass 之间的 partial-PSUM 反馈和下一 pass 准备气泡。对于 $K_{total}>K_{tile}$ 的层，第一轮 K pass 只能产生部分和，后续 pass 必须读取前一轮 partial sum 并继续累加，最后一轮才进入重量化和输出路径。对一个输出 pixel 和一个 $Cout_{tile}$ 通道块，硬件内部形成
\begin{equation}
P[p,c] = \sum_{k=k_0}^{k_0+K_{tile}-1} X[p,k]\cdot W[k,c].
\end{equation}
当当前 pass 不是 final pass 时，该 packet 被写入片上 partial-PSUM buffer；下一 pass 读取同一地址的 partial sum 作为初值继续累加。当当前 pass 是 final pass 时，该 packet 被送入 requant/activation/pool/OFM path。

该语义要求 partial-PSUM 地址、Cout block 和 pass context 与下一 pass 严格对齐，同时避免读写冲突，并保持 final pass 输出顺序与软件重排和 golden 比对一致。本文使用 ping-pong bank 和 per-pass context 维护这些条件。也就是说，PSUM 反馈不是独立于数据流之外的模块，而是跨 K pass 数据流的一部分。

若每轮 pass 都严格串行执行“计算完成、全部 drain、下一轮 feed、下一轮 compute”，后端大层会出现明显阵列空泡。本文首先将 drain writer 改为 AXI-Stream 风格握手：packet valid、addr 和 data 在真实 $valid \& ready$ 握手前保持稳定，握手成功后才推进地址。随后加入 read-return tracker 和 one-entry hold register，使下游 ready 时可以接近每周期输出一个 PSUM packet。

Early drain 允许调度器在当前 pass 已经开始产生 PSUM 数据后提前启动 drain，只要 partial-PSUM FIFO 已具备可读数据。Continuous PSUM collector 进一步改变非 final pass 的路径：collector 在 compute 开始前接收 pass context，并在阵列列 FIFO 产生完整 packet 后立即将其写入 partial-PSUM bank；final pass packet 则进入后处理流水线。这样，非 final pass 的 partial-PSUM 写回从“compute 后集中 drain”变为“compute 过程中持续收集”，减少 pass 边界处的反馈延迟。

K pass 之间除了 PSUM 反馈，还包含下一 pass 权重加载和 IFM replay。Pass-level prefetch 在当前 pass 已经满足安全条件后提前准备下一 pass 的 weight 和 raw-HWC replay。该机制遵循两个约束：一是不能覆盖当前 pass 正在使用的 PE 权重；二是 next-pass IFM replay 只能进入解耦 FIFO 或安全队列，不能破坏当前 pass 的输入顺序。During-compute prefetch 进一步允许在当前 pass compute 进行时启动下一 pass 权重和 raw-HWC replay 的 staging，从而降低 compute stage 中由于下一 pass 准备不足造成的 idle。

图~\ref{fig:psum_prefetch_timeline} 对比了串行 pass 和 OPF-P 重叠 pass 的时间线。

\begin{figure}[htb]
\centering
\includegraphics[width=\textwidth]{\detokenize{img/picture/partial-psum feedback.png}}
\caption{OPF-P 反馈重叠与 pass 级预取时间线}
\label{fig:psum_prefetch_timeline}
\end{figure}

\section{优化策略小结}

表~\ref{tab:strategy_summary} 汇总了 KCS 执行模型与三类优化策略的定位。KCS 负责将不同卷积层统一映射到 LASA；BSD 作用于软硬件协同层，减少 PS/DMA 服务开销；OCRR 作用于数据排布层，缓解 IFM 布局转换和 tile 边界开销；OPF-P 作用于流水线调度层，隐藏 K-pass 反馈和下一 pass 准备气泡。

\begin{table}[htb]
\centering
\caption{KCS 执行模型与三类优化策略}
\label{tab:strategy_summary}
\begin{tabular}{@{}p{0.12\textwidth}p{0.16\textwidth}p{0.26\textwidth}p{0.30\textwidth}@{}}
\toprule
名称 & 层次 & 主要瓶颈 & 对应机制 \\
\midrule
KCS & 执行模型 & 层形状到固定阵列的映射 & K-pass、Cout block、spatial tile \\
BSD & 软硬件协同 & PS 服务与 DMA 启动 & batch stream、packet count、prepacked weight \\
OCRR & 数据排布转换 & IFM 重排与 tile 边界 & raw-HWC cache、replay、backend full-tile \\
OPF-P & 流水线调度 & PSUM/pass 边界气泡 & drainpipe、continuous collector、pass prefetch \\
\bottomrule
\end{tabular}
\end{table}

三类优化共同目标是减少 $T_{noncomp}$，但它们针对的不是同一种开销。这样的分层表述使本文的优化链条从零散实现点收敛为系统协同、数据排布和流水线调度三个层面的策略。下一章将说明这些策略如何通过 runtime、寄存器接口、验证流程和性能计数器集成为可运行、可诊断的板级系统。
